\documentclass[10pt]{article}
\usepackage[margin=0.85in]{geometry}
\linespread{0.97}
\usepackage{amsmath,amssymb,amsthm}
\usepackage[hidelinks]{hyperref}
\setlength{\abovedisplayskip}{5pt plus 1pt minus 1pt}
\setlength{\belowdisplayskip}{5pt plus 1pt minus 1pt}
\setlength{\abovedisplayshortskip}{2pt plus 1pt}
\setlength{\belowdisplayshortskip}{3pt plus 1pt}
\renewcommand{\arraystretch}{1.05}
\newtheorem{theorem}{Theorem}
\newtheorem{proposition}[theorem]{Proposition}
\newtheorem{lemma}[theorem]{Lemma}
\newtheorem{corollary}[theorem]{Corollary}
\theoremstyle{remark}
\newtheorem{remark}[theorem]{Remark}
\newcommand{\DOI}[1]{\href{https://doi.org/#1}{\nolinkurl{#1}}}
\newcommand{\us}{\_\allowbreak}
\emergencystretch=2em
\hyphenation{Zenodo}
\def\UrlBreaks{\do\/}
\def\UrlBigBreaks{}
\newcommand*{\PAPERDOI}{10.5281/zenodo.24407691}
\newcommand*{\SUPPLEMENTDOI}{10.5281/zenodo.24407693}
\newcommand{\kB}{k_{\mathrm B}}
\newcommand{\Tlin}{T_{\mathrm{lin}}}
\newcommand{\Cinf}{C_\infty}
\newcommand{\Ccent}{C_{\mathrm{cent}}}
\newcommand{\Ccaus}{C_{\mathrm{caus}}}
\newcommand{\Cstar}{C_{*}}
\newcommand{\Crep}{C_{\mathrm{rep}}}
\newcommand{\dstar}{d_{*}}
\newcommand{\drep}{d_{\mathrm{rep}}}
\newcommand{\reff}{r_{\mathrm{eff}}}

\hypersetup{
  pdftitle={Centered multichannel Eliashberg branches: a sqrt(341/250)/(2 pi) moment-norm ceiling},
  pdfauthor={Davis Rens},
  pdfsubject={An upper bound on the linear Eliashberg onset temperature of a centered multichannel model},
  pdfkeywords={linearised Eliashberg theory, onset temperature, upper bound, computer-assisted proof}
}
\title{Centered multichannel Eliashberg branches:\\ a $\sqrt{341/250}/(2\pi)$ moment-norm ceiling}
\author{Davis Rens\\SuperMatics Inc.\\\href{https://orcid.org/0000-0001-5839-5197}{ORCID 0000-0001-5839-5197}}
\date{November 2026}

\begin{document}
\maketitle

\begin{abstract}
We prove an upper bound on the linear onset temperature of a reciprocal multichannel Eliashberg model whose channels carry different bounded, energy-dependent, even densities of states and a common retarded glue. The bound reads $\kB\Tlin<C\sqrt{\rho}$, where $\rho$ is the norm of the glue moment matrix, and the constant is $C=\Cstar=\sqrt{341/250}/(2\pi)\approx0.1858777115$ for the supremum over the class and $C=\Crep=\sqrt{6821/5000}/(2\pi)\approx0.1858913384$ as a strict bound for every model of the class. Both numbers lie between the exact scalar constant $\Cinf\approx0.1827262478$ and the constants $\sqrt{683/500}/(2\pi)\approx0.1860139355$ and $\sqrt{41/30}/(2\pi)\approx0.1860593214$ of the preceding preprint, and both are certified by exact rational arithmetic. The relative excess over $\Cinf$ falls from $1.82\%$ to $1.73\%$, and the new constant sits $0.25\%$ above the floor $0.185411$ that the same edge-reserve method cannot go below. The gain comes from a reserve raised from $317/1000$ to $159/500$, a finer partition of the head weights, and a minimisation over the exact rational tangent envelope in the inverse head weight. The finite part of the argument is a covering of a compact rectangle in two variables by closed rational rectangles, carried out twice by independent programs with $159$ and $466$ leaves and guarded by twelve deliberately corrupted runs. The result is an upper bound on a linear model onset temperature. It is not an exact transition temperature and it is not a measurement. The question whether the supremum over the class equals $\Cinf$ remains open.
\end{abstract}

\section{Introduction}

Linearised Eliashberg theory asks when a normal state first becomes unstable to pairing. One fixes the normal branch, takes the pairing interaction to be a retarded glue spectrum, and studies a quadratic form whose negative directions mark the temperatures at which a gap function can grow. A moment ceiling bounds the largest such temperature by a single number built from the glue, its moment norm $\rho$, with the electronic structure entering only through the height and the profile of the density of states. Such a ceiling is a theorem about a model class: it is not a bound on the transition temperature of any material, whose normal state, vertex corrections and competing orders lie outside the model, and nothing below should be read as a prediction about a compound or as a claim that any material is superconducting.

For a single channel the ceiling is settled. For an arbitrary bounded energy-dependent density of states it was proved in~\cite{finite,predecessor} that $\kB\Tlin<\Cinf\sqrt{\rho}$ with $\Cinf=0.18272624777232952$, and that no smaller constant works for that class. Several coupled channels with different, energy-dependent densities of states are harder. The pairing term joining two channels carries the product of their density-of-states averages, while the normal term that opposes pairing in a given channel carries only that channel's own average. A narrow channel therefore weakens its own normal term and its pairing with a wide partner at the same rate, but leaves the wide channel's normal term untouched, so the balance that produces $\Cinf$ in one channel is not available channel by channel; controlling that coupling is the task.

Three earlier preprints gave constants for this coupled problem, for nested classes of the same model and by the same edge-reserve method. The first~\cite{predecessor} obtained $U=\sqrt{6466219/3576050}/(2\pi)\approx0.2140146706$ on centered branches. The second~\cite{optimum} lowered this to $C_1=\sqrt{413/250}/(2\pi)\approx0.2045620904$ by an additive first-frequency comparison, and then to $C_{21}=\sqrt{377/250}/(2\pi)\approx0.1954433070$ by retaining twenty tail frequencies; both hold on the wider causal class of Section~\ref{sec:setting}, which contains the centered one. The third~\cite{finiteenergy} reached $\sqrt{41/30}/(2\pi)\approx0.1860593214$ on centered branches and $\sqrt{45/32}/(2\pi)\approx0.1887345454$ on a wider class of causal branches defined by Harnack inequalities alone, which is $1.82\%$ and $3.29\%$ above $\Cinf$.

The present paper improves the centered constant of~\cite{finiteenergy}, and only that one. We prove
\[
S_{\mathrm{cent}}\ \le\ \Cstar=\frac{\sqrt{341/250}}{2\pi}\approx0.1858777115,
\qquad
\kB\Tlin\ <\ \Crep\sqrt{\rho},\quad \Crep=\frac{\sqrt{6821/5000}}{2\pi}\approx0.1858913384,
\]
where $S_{\mathrm{cent}}$ is the supremum of $\kB\Tlin/\sqrt{\rho}$ over the whole centered class at all finite glue energies. The first is a bound on a supremum and is not strict. The second is strict for every model of the class, and it is the statement one would quote. The relative excess over $\Cinf$ falls from $1.82\%$ to $1.73\%$, so the gain is $0.09\%$; it is small and we do not disguise that. What makes it worth recording is that it is obtained where the method is close to its own ceiling, so the remaining gap to $\Cinf$ is mostly a gap the method cannot cross.

Three changes produce the gain. The reserve $d$, the fraction of the quadratic form each coordinate keeps unconditionally, rises from $317/1000$ to $159/500$, one part in a thousand, and the infinite-tail budget, both cascade stages, the head estimate and the three exterior regions were recomputed and re-certified at the new value. Second, the two head weights are partitioned into forty-three consecutive intervals of a single weight instead of a nine-by-nine grid of pairs. Third, the lower envelope of the centered weights is minimised over the exact convex piecewise-affine envelope in the inverse head weight instead of being evaluated at its worst point, and the programs verify rationally that the breakpoints are ordered.

The last change has a limit that we state at once. The edge-reserve argument of~\cite{finiteenergy} is subject to a floor: an inequality for one edge, valid for every weight sequence of the centered class at every energy and imbalance, cannot produce a constant below $0.185411$. That floor is a statement about the method, not a lower bound for $S_{\mathrm{cent}}$. The constant proved here, $0.1858777115$, lies above it by $0.25\%$, and the floor sits at the fraction $0.852$ of the interval $[\Cinf,\Cstar]$, so the paper closes most of what the method can close.

Whether $S_{\mathrm{cent}}$ equals $\Cinf$ or lies strictly above it is not settled by any of this, and we do not conjecture an answer. What the paper does settle is the comparison at the level of constants: $\Cinf<\Cstar<\Crep<\sqrt{683/500}/(2\pi)<\sqrt{41/30}/(2\pi)<\Ccaus<C_{21}<C_1<U$, a chain of strict inequalities between explicit numbers.

Section~\ref{sec:setting} fixes the model, the branch class and the notion of onset, and states the theorem. Section~\ref{sec:reduction} reduces it to one pair of channels at one energy by a positive-resolvent argument that also explains why the glue moment matrix need not be positive semidefinite. Section~\ref{sec:weights} proves the centered envelope and Section~\ref{sec:proof} the edge inequality case by case. Section~\ref{sec:cert} describes the finite certificate and its limits. Section~\ref{sec:compare} places the constants in a table and states the open question. Section~\ref{sec:repro} explains how to replay everything.

\section{Setting and main result}\label{sec:setting}

\subsection{The model and its normal branch}

Write $\tau=\kB T$. The channel space is a finite set with positive height weights $\eta_x$. Channel $x$ has a measurable density of states $0\le N_x(\xi)\le N_{\max,x}$ whose peak height is the weight of the channel, $N_{\max,x}=\eta_x$, with energy-independent radial vertices. The glue is a family of reciprocal nonnegative measures $d\mathfrak m_{xy}(E)=d\mathfrak m_{yx}(E)\ge0$ on energies $E\in(0,\infty)$, with moment matrix $\mathcal M_{xy}=\int d\mathfrak m_{xy}(E)$. We let $\mathcal M$ act on $\ell^2(\eta)$ by
\[
(\mathcal M h)_x=\sum_y\eta_y\mathcal M_{xy}h_y,\qquad \rho=\|\mathcal M\|_{\ell^2(\eta)}\in(0,\infty).
\]
Entrywise nonnegativity is required and positive semidefiniteness of $\mathcal M$ is not. Since a bipartite glue, one that couples two groups of channels without coupling within a group, has a moment matrix with a negative eigenvalue, this is a genuine relaxation and not a technical convenience. Each measure is supported on positive energies and has finite moment; at each tested temperature its zero-transfer mass is finite.

Pairing is tested on a positive reflected normal branch, that is, on an actual solution of the matched imaginary-axis normal equation. At each tested temperature the frequency rows of the normal kernel, including the zero-transfer term, are assumed to have a finite bound $R(\tau)$; no bound uniform in $\tau$ is needed. The normal kernel is allowed to exceed the pairing transfer by an odd-folded nonnegative excess, a bounded pairing kernel dominated entrywise by the nonnegative attractive envelope cannot raise the onset, and a positive semidefinite static pairing-only subtraction is harmless on the matched branch.

With $\omega_n=\pi\tau w_n$, $w_n=2n+1$, and the normalized profile $\nu_x=N_x/N_{\max,x}$ even and valued in $[0,1]$, the weights are the Poisson averages
\begin{equation}\label{eq:weights}
b_{n,x}=\int_{-\infty}^{\infty}\nu_x(\xi)\,\frac{\pi\tau s_{n,x}}{(\pi\tau s_{n,x})^2+\xi^2}\,\frac{d\xi}{\pi}\in[0,1],
\end{equation}
where the normal width at the $n$th reflected frequency, in units in which the free width is $w_n$, is
\begin{equation}\label{eq:normal}
s_{n,x}=w_n+\frac{1}{(2\pi\tau)^2}\sum_y\eta_y\int\sum_{m\ge0}p_{nm}(t)\,b_{m,y}\,d\mathfrak m_{xy}(E)\ \ge\ w_n ,
\end{equation}
with $t$ and $p_{nm}$ defined below. An extra allowed normal mass only increases the displayed width lower bound, so nothing below depends on assuming the minimal normal kernel.

The width cone of the centered class is
\begin{equation}\label{eq:cone}
w_m^{-1}\ \le\ \frac{s_{m,x}}{s_{0,x}}\ \le\ w_m,\qquad m\ge1,
\end{equation}
which follows from $s_{n,x}/w_n$ nonincreasing and $w_ns_{n,x}$ nondecreasing in $n$; it is established for the centered class in~\cite{predecessor,optimum}. Because the two centered Poisson kernels at heights whose ratio lies in $[w_m^{-1},w_m]$ have a pointwise ratio in the same interval, integrating both $\nu_x$ and its complement gives the four inequalities
\begin{equation}\label{eq:harnack}
0\le b_{m,x}\le1,\qquad \frac{b_{0,x}}{w_m}\le b_{m,x}\le w_mb_{0,x},\qquad
w_mb_{0,x}-(w_m-1)\le b_{m,x}\le\frac{b_{0,x}+w_m-1}{w_m}.
\end{equation}
We use the centered class throughout, so \eqref{eq:cone} and the sharper lower envelope of Section~\ref{sec:weights} are both available. The causal class of~\cite{finiteenergy}, in which only \eqref{eq:harnack} is retained, is not treated here and its constant is not improved.

\subsection{The stability form and the onset}

For a spectral energy $E$ put
\[
t=\Bigl(\frac{E}{2\pi\tau}\Bigr)^2,\qquad k_j(t)=\frac1{t+j^2},\qquad p_{nm}=k_{|n-m|}-k_{n+m+1}\ge0,\qquad S_{nm}=k_{|n-m|}+k_{n+m+1},
\]
the energy-squared factor removed from the physical transfer being $(2\pi\tau)^2$. For a finite-support real profile $v_{n,x}$ put $g_{n,x}=\sqrt{w_nb_{n,x}}\,v_{n,x}$. After folding the frequency axis the stability form is
\begin{equation}\label{eq:form}
Q(v)=\sum_x\eta_x\|g_x\|^2+\frac{1}{2(2\pi\tau)^2}\sum_{x,y}\eta_x\eta_y\int I_{xy}(t)\,d\mathfrak m_{xy}(E),
\end{equation}
where, with $b_n=b_{n,x}$, $c_m=b_{m,y}$, $X_n=v_{n,x}$ and $Y_m=v_{m,y}$,
\[
I_{xy}(t)=\sum_{n,m\ge0}b_nc_m\bigl[p_{nm}(X_n^2+Y_m^2)-2S_{nm}X_nY_m\bigr].
\]
The onset $\Tlin$ is the supremum of the temperatures at which $Q$ fails to be positive definite on finitely supported real profiles of a single symmetry sector; for a model with several sectors it is the largest of the sector onsets. It is an energy-valued linear onset temperature. It is not the transition temperature of an exact Hamiltonian, and it is not measured.

\subsection{The theorem}

\begin{theorem}\label{thm:main}
In the centered class above, with $\rho$ the norm of the glue moment matrix,
\[
S_{\mathrm{cent}}:=\sup\frac{\kB\Tlin}{\sqrt{\rho}}\ \le\ \Cstar=\frac{\sqrt{341/250}}{2\pi},
\]
the supremum being over all models of the class with finitely many channels, all finite positive glue energies, and all admissible centered branches; and for every such model
\[
\kB\Tlin\ <\ \Crep\sqrt{\rho},\qquad \Crep=\frac{\sqrt{6821/5000}}{2\pi}.
\]
The two constants satisfy the exact rational decimal enclosures
\[
0.185877711471916<\Cstar<0.185877711471917,\qquad
0.185891338370617<\Crep<0.185891338370618 .
\]
\end{theorem}

The two constants carry the two reserves $\dstar=159/500=0.318$ and $\drep=3179/10000=0.3179$, through $2(1-\dstar)=341/250$ and $2(1-\drep)=6821/5000$. The strict bound is not a by-product of a limit. It follows from $S_{\mathrm{cent}}\le\Cstar$ with $\Cstar<\Crep$, and at $\tau=\Crep\sqrt\rho$ the bare gap of the form equals
\begin{equation}\label{eq:baregap}
1-\frac{1-\dstar}{1-\drep}=\frac1{6821}>0 ,
\end{equation}
so that temperature is itself stable by a rational margin. The strict inequality at each model is thus a rational statement about that model, and it does not require any assertion that stability at one temperature forces stability in a neighbourhood of it. No monotonicity of the onset in temperature is used anywhere.

\begin{remark}\label{rem:notsharp}
Neither constant is claimed exact. The exact scalar constant is $\Cinf\approx0.1827262478$ and $\Cinf\le S_{\mathrm{cent}}$ holds because the scalar case is a member of the class and its supremum over the class is exactly $\Cinf\sqrt\rho$; see~\cite{finite,predecessor}. Thus $\Cinf\le S_{\mathrm{cent}}\le\Cstar$ with $\Cstar/\Cinf-1=1.725\%$. Whether the lower end or a point strictly above it is correct is open; Section~\ref{sec:compare} states what is and is not known.
\end{remark}

\section{Reduction to an edge inequality}\label{sec:reduction}

Real and imaginary parts of a trial function decouple in the form \eqref{eq:form}, so real trials suffice. For $r>0$ put
\[
B_r=\sum_{n\ge0}w_n\bigl(r^{-1}b_nX_n^2+r\,c_nY_n^2\bigr)=r^{-1}\|g_x\|^2+r\|g_y\|^2 .
\]

\begin{proposition}[edge inequality]\label{prop:edge}
With $\dstar=159/500$, for every $t>0$, every $r>0$, all centered-admissible weight sequences $(b_n)$ and $(c_m)$, and all real finitely supported $X,Y$,
\begin{equation}\label{eq:edge}
B_r+\tfrac12 I_{xy}(t)\ \ge\ \dstar\,B_r .
\end{equation}
\end{proposition}

The two weight sequences are independent, so the diagonal case $x=y$ is included, and the constant $\dstar$ carries no dependence on $t$, on $r$ or on the weights.

\begin{proof}[Proof of Theorem~\ref{thm:main} from Proposition~\ref{prop:edge}]
Choose $R'>\rho$ and put $h=(I-\mathcal M/R')^{-1}\mathbf 1=\sum_{j\ge0}(\mathcal M/R')^j\mathbf 1$. Every entry of $\mathcal M$ is nonnegative and $\mathbf 1>0$, so $h>0$ entrywise and $\mathcal Mh=R'(h-\mathbf 1)\le R'h$ entrywise; no positivity of the quadratic form of $\mathcal M$ is used. Apply \eqref{eq:edge} with $r=r_{xy}=h_x/h_y$ for every ordered pair and integrate against $d\mathfrak m_{xy}$. Then
\[
Q(v)\ \ge\ \sum_x\eta_x\|g_x\|^2-\frac{1-\dstar}{(2\pi\tau)^2}\sum_{x,y}\eta_x\eta_y\mathcal M_{xy}\Bigl(\frac{h_y}{h_x}\|g_x\|^2+\frac{h_x}{h_y}\|g_y\|^2\Bigr).
\]
By reciprocity $\mathcal M_{xy}=\mathcal M_{yx}$ the two terms in the bracket give the same double sum, and $\sum_y\mathcal M_{xy}\eta_yh_y\le R'h_x$ bounds it by $R'\sum_x\eta_x\|g_x\|^2$. Hence
\begin{equation}\label{eq:gap}
Q(v)\ \ge\ \Bigl[1-\frac{2(1-\dstar)R'}{(2\pi\tau)^2}\Bigr]\sum_x\eta_x\|g_x\|^2 ,
\end{equation}
and the left side does not depend on $R'$, so we may let $R'\downarrow\rho$.

Now fix $\tau>\Cstar\sqrt\rho$, so that $(2\pi\tau)^2>2(1-\dstar)\rho$, and put $\beta(\tau)=1-2(1-\dstar)\rho/(2\pi\tau)^2>0$. Letting $R'\downarrow\rho$ in \eqref{eq:gap} gives $Q(v)\ge\beta(\tau)\sum_x\eta_x\|g_x\|^2$ for every $v$. The normal widths obey $s_{n,x}\le w_n+R(\tau)\le(1+R(\tau))w_n$ at this temperature, since every frequency row of the normal kernel is bounded by $R(\tau)$, so the Hilbert norm of the pairing problem, $\sum_x\eta_x\|(s_xb_x)^{1/2}v_x\|^2$, is at most $(1+R(\tau))$ times $\sum_x\eta_x\|g_x\|^2$. Therefore
\[
Q(v)\ \ge\ \frac{\beta(\tau)}{1+R(\tau)}\sum_x\eta_x\|(s_xb_x)^{1/2}v_x\|^2
\]
for every $v$, and $Q$ is strictly positive definite at $\tau$. No $\tau>\Cstar\sqrt\rho$ is a failure temperature, so $\kB\Tlin\le\Cstar\sqrt\rho$; disconnected instability sets are included, since the argument is applied at each $\tau$ and not along a branch.

For the strict statement, take $\tau\ge\Crep\sqrt\rho$ with $(2\pi\tau)^2\ge2(1-\drep)\rho$. Then the bracket in \eqref{eq:gap} is at least $1-(1-\dstar)R'/((1-\drep)\rho)$, and letting $R'\downarrow\rho$ gives
\[
Q(v)\ \ge\ \frac1{6821}\sum_x\eta_x\|g_x\|^2 .
\]
The bare gap at the reported temperature is positive and its stability operator is bounded below by $[6821(1+R(\tau))]^{-1}>0$, so that temperature is itself stable, and $\Cstar<\Crep$ turns $\kB\Tlin\le\Cstar\sqrt\rho$ into the strict statement $\kB\Tlin<\Crep\sqrt\rho$ quoted in the theorem. Both constants are functions of the raw reserves through $2(1-d)$ and of the reported deficit through the displayed ratio, and the identity \eqref{eq:baregap} is a rational identity.
\end{proof}

All the work therefore lies in Proposition~\ref{prop:edge}. Note that the previous preprint used $\dstar=317/1000$, which gives $2(1-\dstar)=683/500$; the single number that changes is the reserve.

\section{The centered weight envelope}\label{sec:weights}

The completion of Section~\ref{sec:proof} needs lower bounds for the tail weights in terms of the head weight, uniform over the centered class. They come from the Poisson representation \eqref{eq:weights} and the cone \eqref{eq:cone}, and they are sharper than the Harnack lines \eqref{eq:harnack}.

\begin{proposition}[Poisson envelope]\label{prop:poisson}
Let $w>1$, let $N:[0,\infty)\to[0,1]$ be measurable, and put
\[
b_0=\frac2\pi\int_0^\infty\frac{N(\xi)}{1+\xi^2}\,d\xi,\qquad
b(\sigma)=\frac2\pi\int_0^\infty N(\xi)\,\frac{\sigma}{\sigma^2+\xi^2}\,d\xi .
\]
Then $b(\sigma)\ge\varphi_w(b_0)$ for every $\sigma\in[w^{-1},w]$, where
\begin{equation}\label{eq:phi}
\varphi_w(b)=\frac2\pi\arctan\Bigl(\frac1w\tan\frac{\pi b}2\Bigr).
\end{equation}
\end{proposition}

\begin{proof}
If $b_0\in\{0,1\}$ then $N$ vanishes or equals $1$ almost everywhere and the claim is immediate. Otherwise put $W=\tan(\pi b_0/2)\in(0,\infty)$. We minimise $\int Nk_\sigma$ with $k_\sigma(\xi)=\sigma/(\sigma^2+\xi^2)$ over measurable $N$ with values in $[0,1]$ and a prescribed value of $\int Nk_1$. By the bathtub principle, a form of the Neyman-Pearson lemma, a minimiser is the indicator of a level set of the density ratio
\[
\frac{k_\sigma}{k_1}=\frac{\sigma(1+\xi^2)}{\sigma^2+\xi^2},
\]
whose derivative in $x=\xi^2$ has the sign of $\sigma^2-1$. For $\sigma\in(1,w]$ the ratio increases with $\xi$, so the minimiser fills an interval adjacent to zero, namely $\mathbf 1_{[0,W]}$, which has the prescribed $b_0$, and
\[
b(\sigma)\ \ge\ \frac2\pi\arctan\frac W\sigma\ \ge\ \frac2\pi\arctan\frac Ww=\varphi_w(b_0).
\]
For $\sigma\in[w^{-1},1)$ the ratio decreases, so the minimiser fills a tail $\mathbf 1_{[1/W,\infty)}$, which also has the prescribed $b_0$, and $b(\sigma)\ge\frac2\pi\arctan(\sigma W)\ge\varphi_w(b_0)$. For $\sigma=1$, $b(\sigma)=b_0\ge\varphi_w(b_0)$. Arbitrary measurable $N$ follow by exchanging mass between points ordered by the ratio, which is the content of the bathtub principle, and the endpoint cases by continuity.
\end{proof}

In the application the energy variable is rescaled by $\pi\tau s_{0,x}$, so $N$ is the normalised density $\nu_x$, $b_0=b_{0,x}$ and $b_m=b(\sigma)$ with $\sigma=s_{m,x}/s_{0,x}\in[w_m^{-1},w_m]$ by \eqref{eq:cone}.

For $w>1$ the function $\varphi_w$ is increasing and convex on $[0,1]$. With $\xi=\tan(\pi b/2)$,
\[
\varphi_w'(b)=\frac{w(1+\xi^2)}{w^2+\xi^2},\qquad
\varphi_w''(b)=\frac{\pi w(w^2-1)\,\xi(1+\xi^2)}{(w^2+\xi^2)^2}>0 ,
\]
so every tangent line lies below $\varphi_w$ and gives an admissible lower bound,
\begin{equation}\label{eq:tangent}
b_m\ \ge\ L_{\xi,m}\,b_0+C_{\xi,m},\qquad
L_{\xi,m}=\frac{w_m(1+\xi^2)}{w_m^2+\xi^2},\qquad
C_{\xi,m}=\frac2\pi\arctan\frac{\xi}{w_m}-L_{\xi,m}\frac2\pi\arctan\xi .
\end{equation}
We use the reference points $\xi\in\{4,5,6,8,10,12,16,24\}$. The intercept is not positive, because $\varphi_w(0)=0$ and $\varphi_w$ is convex, and the bound \eqref{eq:tangent} is valid with a negative intercept as well; only coefficients of nonnegative sign may multiply such a bound. Together with the two Harnack lines $b_m\ge b_0/w_m$ and $b_m\ge w_mb_0-(w_m-1)$, the ten lines make the centered envelope a maximum of ten affine functions of $b_0$,
\begin{equation}\label{eq:env}
b_m\ \ge\ \max_{\ell}\bigl(\lambda_{\ell,m}b_0+\nu_{\ell,m}\bigr),\qquad \lambda_{\ell,m}\ge\frac1{w_m} .
\end{equation}
The intercepts \eqref{eq:tangent} are irrational. The rational numbers used are obtained from certified alternating-series enclosures of $\pi$ and of the arctangent in integer arithmetic, with every lower bound rounded downward, and the two implementations use different series lengths and separate arctangent routines.

\section{Proof of the edge inequality}\label{sec:proof}

\subsection{Coordinates and four disjoint sectors}

Fix $t>0$, $r>0$ and two admissible weight sequences $b_n$ and $c_m$. Put
\[
a_n=r^{-1/2}\sqrt{w_nb_n}\,X_n,\qquad z_m=r^{1/2}\sqrt{w_mc_m}\,Y_m,
\]
so that $B_r=\sum_{n\ge0}(a_n^2+z_n^2)$. If a weight vanishes the corresponding terms drop out, and the bare reserve $1\ge\dstar$ is untouched; zero-weight coordinates are removed by continuity. The part of $\tfrac12I_{xy}$ that couples frequency $n$ of the first channel with frequency $m$ of the second is
\begin{equation}\label{eq:block}
\mathcal B_{nm}=\frac{r\,c_m\,p_{nm}}{2w_n}\,a_n^2+\frac{b_n\,p_{nm}}{2r\,w_m}\,z_m^2-S_{nm}\sqrt{\frac{b_nc_m}{w_nw_m}}\;a_nz_m .
\end{equation}
The unit terms $a_n^2$ and $z_n^2$ of $B_r$ are the mass we spend. Of each we keep back the fraction $\dstar$ demanded by \eqref{eq:edge}, and the remaining $1-\dstar$ is shared among the pairs in which the coordinate takes part.

Fix $M=20$ and split the set of all frequency pairs into four parts, which are disjoint and cover it:
\begin{enumerate}
\item the head pair $(0,0)$;
\item the head-tail pairs $(0,m)$ and $(n,0)$ with $m,n\ge1$, called Stage~0;
\item the pairs $(1,k)$ and $(k,1)$ with $2\le k\le M$, called Stage~1;
\item all remaining tail-tail pairs with $n,m\ge1$.
\end{enumerate}
Every tail-tail pair beyond the chosen ones is bounded below by its row cost, so no pair is left unaccounted.

\subsection{Reflected tails and the infinite tail reduction}

The identity
\[
(a-b)(X^2+Y^2)-2(a+b)XY=(a-b)(X-Y)^2-4bXY
\]
with $a=k_{|n-m|}$ and $b=k_{n+m+1}$ writes the pair $(n,m)$ of \eqref{eq:form} as a nonnegative difference square plus the reflected cross term $-2b_nc_mk_{n+m+1}X_nY_m$. For the remaining tail-tail pairs we discard the difference square and bound the cross term by the weighted arithmetic-geometric inequality,
\[
-2b_nc_mk_{n+m+1}X_nY_m\ \ge\ -\frac{a_n^2}{w_m(n+m+1)^2}-\frac{z_m^2}{w_n(n+m+1)^2},
\]
using $b_n,c_m\le1$ and $k_{n+m+1}\le(n+m+1)^{-2}$. The factor $2$ in the reflected term is what makes this step valid, and dropping it would invalidate the bound. Collecting the coefficient of $a_n^2$ over the remaining pairs gives, for $1\le n\le M$,
\begin{equation}\label{eq:qred}
q_n=\sum_{\substack{1\le m\le M\\(n,m)\notin\{(1,k),(k,1):\,2\le k\le M\}}}\frac1{w_m(n+m+1)^2}+\frac1{43(n+M+1)},
\end{equation}
where the last term bounds the modes $m\ge M+1$, for which $w_m\ge43$ and
\[
\sum_{m\ge M+1}\frac1{(n+m+1)^2}\ \le\ \int_{M+1}^{\infty}\frac{dm}{(n+m)^2}=\frac1{n+M+1}.
\]
For $n>M$ no Stage-1 pair is removed, so the full row sum applies, and it is at most that of the first row: the terms decrease in $n$ and the tail term does too. Evaluating sixteen terms exactly and bounding the rest gives
\begin{equation}\label{eq:qcap}
q_{\mathrm{full},n}\ \le\ \sum_{m=1}^{16}\frac1{w_m(m+2)^2}+\frac1{4\cdot16^2}\ <\ \frac{33}{500}.
\end{equation}
These are rational infinite-tail bounds. They are not a truncation of the physical normal equation: the physical tail is bounded, not cut.

\subsection{The cascade}

After the reserve and these costs, tail mode $n\le M$ has the share $\mathrm{rem}_n=1-\dstar-q_n$ of its bare mass left, and $\mathrm{rem}_1=1-159/500-973/25542\approx0.6439$. We give $\mathrm{rem}_1$ entirely to Stage~0 and split $\mathrm{rem}_k$ for $2\le k\le M$ in the proportion $7/10$ to Stage~0 and $3/10$ to Stage~1,
\[
\alpha_{0,1}=\mathrm{rem}_1,\quad \beta_{1,1}=0;\qquad
\alpha_{0,k}=\tfrac7{10}\mathrm{rem}_k,\quad \beta_{1,k}=\tfrac3{10}\mathrm{rem}_k\quad(2\le k\le M),
\]
and $\alpha_{\rm ext}=1-\dstar-33/500=77/125$ for $n>M$. All remainders are positive, and in the range where Stage~1 is used one has $\beta_{1,k}>0.1984$.

\emph{Stage 1.} For $n=1$ and $2\le k\le M$ the share is $\beta_{1,k}$, and $b_1$ drops out in the worst case $b_1=1$, because $b_1\mapsto b_1/(\beta+Cb_1)$ is increasing in $b_1$. Completing the square in $z_k$ against the coordinate $a_1$ leaves $c_kD_{1k}$ in $a_1^2$, with
\begin{equation}\label{eq:D1}
D_{1k}=\frac{r\,p_{1k}}6-\frac{S_{1k}^2}{12w_k\bigl(\beta_{1,k}+r^{-1}p_{1k}/(2w_k)\bigr)}
=\frac{2p_{1k}\bigl(rw_k\beta_{1,k}-2/(3w_k)\bigr)}{12w_k\beta_{1,k}+6r^{-1}p_{1k}} ,
\end{equation}
the second form following from the exact identity $S_{1k}^2-p_{1k}^2=4p_{1k}/(3w_k)$, which holds because $p_{1k}=k_{k-1}-k_{k+2}$ and $(k+2)^2-(k-1)^2=3w_k$. Its numerator is positive for $r\in[1/3,3]$ and $k\ge2$, so no sign case arises; on that range it increases in $r$ and in $p_{1k}$ and therefore decreases in $t$. The first tail mode of the first channel is thereby boosted to
\[
\alpha_{0,1}^{\rm eff}=\alpha_{0,1}+\sum_{k=2}^{M}c_kD_{1k},
\]
and the mirror statement holds for the first tail mode of the second channel, with $r$ replaced by $r^{-1}$ and $b$ by $c$: $D^{Y}_{1k}(r,t)=D_{1k}(r^{-1},t)$. This is the cascade. The first tail mode is stabilised by the higher tail modes before it meets the head, and the boost is what allows the reserve to exceed the value at which the single-stage completion of the first mode turns negative.

\emph{Stage 0.} The bare share of the tail coordinate $z_m$ is $\alpha_m$, where $\alpha_1=\alpha_{0,1}^{\rm eff}$ for the second channel is the boosted value and $\alpha_m=\alpha_{0,m}$ for $2\le m\le M$. In the normalisation in which the share is measured against the head weight, this is $A_m=\alpha_m/b_0$, so completing the square in $z_m$ against the head coordinate $a_0$ leaves $c_mD_{0m}$ in $a_0^2$, where
\begin{equation}\label{eq:D0}
D_{0m}=\frac{r\,p_{0m}}2-\frac{S_{0m}^2}{4w_m\bigl(A_m+r^{-1}p_{0m}/(2w_m)\bigr)}
=\frac{2p_{0m}\bigl(rw_mA_m-2/w_m\bigr)}{4w_mA_m+2r^{-1}p_{0m}} ,
\end{equation}
using $S_{0m}^2-p_{0m}^2=4p_{0m}/w_m$. The division by $b_0$ is exactly the dependence on the head weight. A positive numerator term increases in $r$, in $p_{0m}$ and in $A_m$, so it is bounded below using lower values of all three, $p_{0m}$ at the larger $t$ endpoint, and an upper bound for $r^{-1}$ in the denominator. A negative numerator term needs care with the signs: it is bounded below by dropping the positive denominator contribution $2r^{-1}p_{0m}$ and using $p_{0m}$ at the smaller $t$ endpoint, and then by multiplying by the upper Harnack ratio $c_m/c_0\le w_m$, with $b_m/b_0\le w_m$ on the mirror side. Using a lower ratio on a negative term would have the wrong sign, and this is the step where the upper line of \eqref{eq:harnack} is indispensable.

For $m>M$ only $\alpha_{\rm ext}$ is available. Since $w_m\ge43$ and $(43/3)\alpha_{\rm ext}>2/43$, every unretained Stage-0 completion has a positive numerator and may be discarded rather than bounded below. Thus no finite-frequency remainder of unknown sign is left anywhere.

\subsection{The head criterion}

After the two stages the two head coordinates $a_0$ and $z_0$ remain, with budget $\Phi_X$ and $\Phi_Y$. Writing $\tilde a=\sqrt{c_0}\,a_0$ and $\tilde z=\sqrt{b_0}\,z_0$, the head part is the form with matrix
\[
M_{\rm head}=\begin{pmatrix}\Phi_X+\beta r&-(\beta+u)\\-(\beta+u)&\Phi_Y+\beta/r\end{pmatrix},
\qquad \beta=\frac1{2t(1+t)},\quad u=\frac1{1+t},
\]
where
\[
\Phi_X=\frac{1-\dstar}{c_0}+\sum_{m=1}^{M}\frac{c_m}{c_0}D^{X}_{0m},\qquad
\Phi_Y=\frac{1-\dstar}{b_0}+\sum_{m=1}^{M}\frac{b_m}{b_0}D^{Y}_{0m},
\]
and $D^{X}_{0m}$ and $D^{Y}_{0m}$ denote the lower bounds just described. If $b_0=0$ or $c_0=0$ the cross term vanishes and the claim is immediate. For the head pair, $p_{00}=1/(t(1+t))=2\beta$ and $S_{00}=(1+2t)/(t(1+t))=2(\beta+u)$.

\begin{lemma}[head criterion]\label{lem:head}
Let $\Phi_X,\Phi_Y>0$ and $A=2t(1+t)$. Then
\[
A^2r\det M_{\rm head}=(A\Phi_X+r)(Ar\Phi_Y+1)-r(1+2t)^2=A^2r\Phi_X\Phi_Y+A(\Phi_X+r^2\Phi_Y)-4rt(1+t),
\]
and this equals
\[
2rt(1+t)\Bigl[2t(1+t)\Phi_X\Phi_Y+\frac{\Phi_X}{r}+r\Phi_Y-2\Bigr].
\]
Consequently $M_{\rm head}\succeq0$ if and only if
\begin{equation}\label{eq:det}
2t(1+t)\,\Phi_X\Phi_Y+\frac{\Phi_X}{r}+r\Phi_Y\ \ge\ 2 .
\end{equation}
\end{lemma}

\begin{proof}
With $\beta=1/A$ one has $A(\Phi_X+\beta r)=A\Phi_X+r$ and $Ar(\Phi_Y+\beta/r)=Ar\Phi_Y+1$, while $A^2r(\beta+u)^2=r(1+2t)^2$ because $A/(1+t)=2t$. This gives the first equality. Expanding and using $(1+2t)^2=1+4t(1+t)$ gives the second, in which $-4rt(1+t)=-2rA$ combines with the remaining terms to give the factored form. Both diagonal entries are positive, so $M_{\rm head}\succeq0$ if and only if its determinant is nonnegative.
\end{proof}

The criterion \eqref{eq:det} is continuous at $t=0$ after the positive denominators are cleared, so the same expression closes rectangles that touch $t=0$; physical spectral atoms have $t>0$. On a rectangle $[r_{\rm lo},r_{\rm hi}]\times[t_{\rm lo},t_{\rm hi}]$, with $\Phi_X\ge P_X>0$ and $\Phi_Y\ge P_Y>0$, the function $r\mapsto P_X/r+rP_Y$ is convex, so it lies above its tangent line at the midpoint $r_0=\tfrac12(r_{\rm lo}+r_{\rm hi})$,
\[
\frac{P_X}{r}+rP_Y\ \ge\ \frac{2P_X}{r_0}+r\Bigl(P_Y-\frac{P_X}{r_0^2}\Bigr).
\]
The right side is affine in $r$, so its minimum over the rectangle is at $r_{\rm lo}$ when the slope is nonnegative and at $r_{\rm hi}$ otherwise; call it $g_{\rm lo}$. It is increasing in $P_X$ and $P_Y$ on the rectangle because $r<2r_0$, so lower bounds may be inserted. A sufficient condition for \eqref{eq:det} on the whole rectangle is then
\begin{equation}\label{eq:rect}
2t_{\rm lo}(1+t_{\rm lo})\,P_XP_Y+g_{\rm lo}\ \ge\ 2 ,
\end{equation}
a finite rational inequality. All acceptance predicates are of this form, for closed rectangles and closed head intervals, so the covering of Section~\ref{sec:cert} certifies the continuum and not a grid of parameter points.

\subsection{Head weights and their minimisation}

The quantities $\Phi_X$ and $\Phi_Y$ involve the head weights through the ratios $b_m/b_0$ and $c_m/c_0$, which \eqref{eq:env} bounds below by a maximum of affine functions of the head weight; on a head interval $[b_{\rm lo},b_{\rm hi}]$ one takes the smallest value of each. Writing $u_0=1/b_0$, a reserve $A/b_0+\sum_mC_mb_m/b_0$ with $C_m\ge0$ becomes, after \eqref{eq:env}, a maximum of affine functions of $u_0$, hence convex and piecewise affine in $u_0$, so its minimum over the interval is attained at an endpoint or at a crossing of two lines. Both programs verify rationally that the adjacent breakpoints are ordered, for every $m$, so enumerating adjacent crossings omits no switch of the envelope. On the first interval $(0,1/2]$ both verify that every additional line lies at or below the basic line $b_0/w_m$ at $b_0=1/2$, and since each such line has slope at least $1/w_m$ it stays below that line on all of $(0,1/2]$. Hence the envelope ratio is the constant $1/w_m$ there and $A/b_0+\sum C_m/w_m$ is minimised at $b_0=1/2$, so no limit $b_0\to0^+$ is missed.

We partition $[0,1]$ at the knots
\[
0,\ \tfrac12,\ \tfrac34,\ \tfrac{160}{200},\ \tfrac{161}{200},\ \dots,\ \tfrac{200}{200},
\]
which gives $43$ consecutive intervals: the three intervals $(0,\tfrac12]$, $[\tfrac12,\tfrac34]$, $[\tfrac34,\tfrac45]$ and forty intervals of length $1/200$ from $4/5$ to $1$. The coarse part needs the comparison just made at $b_0=1/2$; the fine upper part is where the envelope actually bends. The partition is verified to be gapless with the correct endpoints, so the head weights are covered exactly once.

\subsection{The exterior regions}

The compact rectangle $[1/3,3]\times[0,1/2]$ is where the cascade and the head minimisation are needed. Outside it a cruder accounting suffices, and it uses no Stage~1. Every tail coordinate keeps the share $\alpha_{\rm ext}=77/125$ of its bare mass, because its row cost is at most $33/500$. A nonnegative Stage-0 completion is dropped, and a negative one is charged against the reserve. What remains to be shown is that the two head coordinates keep enough, and for that two unassisted estimates of the head pair are recorded in~\cite{finiteenergy} and reused here:
\begin{equation}\label{eq:headest}
\text{kept fraction of the head pair}\ \ge\ 1-\frac{2}{r+r^{-1}},
\qquad\text{and}\qquad \text{kept fraction}\ \ge\ \frac{t}{1+t}.
\end{equation}
The second is the immediate one: dropping the positive square $\beta(\sqrt r\,\tilde a-\tilde z/\sqrt r)^2$ leaves the factor $1-u=t/(1+t)$ with $u=1/(1+t)$. The first does not degrade as $t\downarrow0$, which is why the far-imbalance case uses it, while the second is what controls large energy. Either estimate, less the charged completions, must exceed $\dstar$; when it does, both head coordinates keep at least the fraction $\dstar$ of their bare mass, and \eqref{eq:edge} follows. Three cases cover the complement of the compact rectangle.

\emph{Far imbalance, $r\le1/4$.} The case $r\ge4$ follows by the channel swap. Here the first estimate in \eqref{eq:headest} is at least $9/17$, since $r+r^{-1}\ge17/4$, and it is uniform in $t$. For every $m\ge1$ and $t\ge0$, dropping the nonnegative term $rw_mA_m$ from the numerator and the term $2r^{-1}p_{0m}$ from the denominator in \eqref{eq:D0} gives
\[
D_{0m}\ \ge\ -\frac{p_{0m}(t)}{w_m^2\alpha_{\rm ext}}\ \ge\ -\frac1{\alpha_{\rm ext}(2m+1)m^2(m+1)^2},
\]
using $p_{0m}\le(2m+1)/(m^2(m+1)^2)$. Since $(2m+1)m^2(m+1)^2>2m^5$, the terms with $m\ge21$ sum to less than $1/(8\cdot20^4)$, so the total charge of the negative completions is less than $(7/75)/\alpha_{\rm ext}$, and
\[
\sum_{m\ge1}\frac1{(2m+1)m^2(m+1)^2}\ \le\ \sum_{m=1}^{20}\frac1{(2m+1)m^2(m+1)^2}+\frac1{8\cdot20^4}\approx0.0903550<\frac7{75}.
\]
The fraction that remains is therefore at least
\[
\frac9{17}-\frac{7/75}{\alpha_{\rm ext}}=\frac{212}{561}\approx0.377897>\dstar .
\]

\emph{Intermediate imbalance, $1/4\le r\le1/3$.} The case $3\le r\le4$ again follows by the swap. Now $r+r^{-1}\ge10/3$, so the first estimate in \eqref{eq:headest} is at least $2/5$. For $m\ge2$ one has $w_m\ge5$ and $r\alpha_{\rm ext}/2\ge77/1000>1/25\ge w_m^{-2}$, so those completions are positive and only $m=1$ can be negative. Using $p_{0,1}\le3/4$ and dropping the positive denominator term, $D_{0,1}\ge\frac34(\alpha_{\rm ext}/8-1/9)/\alpha_{\rm ext}$, and the fraction that remains is at least
\[
\frac25+\frac34\cdot\frac{\alpha_{\rm ext}/8-1/9}{\alpha_{\rm ext}}=\frac{13249}{36960}\approx0.358469>\dstar .
\]

\emph{Large energy, $t\ge1/2$ on $[1/3,3]$.} Here $r$ and $r^{-1}$ both lie in $[1/3,3]$ and the second estimate in \eqref{eq:headest} gives at least $1/3$. The completions with $m\ge2$ are again positive. For $m=1$ one has $p_{0,1}=3/((t+1)(t+4))\le4/9$ and $\alpha_{\rm ext}/6-1/9<0$, so $D_{0,1}\ge\frac49(\alpha_{\rm ext}/6-1/9)/\alpha_{\rm ext}$, and the fraction that remains is at least
\[
\frac13+\frac49\cdot\frac{\alpha_{\rm ext}/6-1/9}{\alpha_{\rm ext}}=\frac{2041}{6237}\approx0.327241>\dstar .
\]
Together with the compact rectangle and its reflection, these three cases exhaust $r>0$, $t>0$.

\begin{proof}[Completion of the proof of Proposition~\ref{prop:edge}]
Fix $t>0$, $r>0$ and admissible $b,c$. If $r\notin[1/3,3]$ or $t\ge1/2$ the exterior cases just proved give both head coordinates and every tail coordinate at least the fraction $\dstar$ of their bare mass, which is \eqref{eq:edge}; adding back the nonnegative squares discarded along the way returns \eqref{eq:block}, and summing the blocks recovers $B_r+\tfrac12I_{xy}\ge\dstar B_r$. If $r\in[1/3,3]$ and $t\in(0,1/2)$, the head intervals containing $b_0$ and $c_0$ are two of the $43$ intervals, and the finite certificate of Section~\ref{sec:cert} shows that the sufficient condition \eqref{eq:rect} holds for that rectangle and that pair of head intervals, with the lower values $P_X,P_Y$ obtained from \eqref{eq:env} and \eqref{eq:D0} and \eqref{eq:D1}. Then \eqref{eq:det} holds, and the same back-substitution gives the edge inequality.
\end{proof}

\section{The finite certificate}\label{sec:cert}

It remains to show that \eqref{eq:rect} holds on the compact domain, for every rectangle of a complete covering and every pair of head intervals, with the reserve $\dstar=159/500$. This is a finite computation in exact rational arithmetic, and it was carried out twice by programs that do not share code.

\emph{The domain and the covering.} The primary program uses the symmetry $r\leftrightarrow r^{-1}$, which exchanges the two channels of a pair and maps the problem to itself, to reduce to $[1,3]\times[0,1/2]$, an area of $1$. It subdivides by bisection, halving the $t$ interval when the $r$ interval is less than three halves of it and the $r$ interval otherwise, and it stops when \eqref{eq:rect} holds at the certified reserve and for all $43^2=1849$ pairs of head intervals. It terminates at depth $20$ with $159$ leaves, and its exact rational leaf areas sum to exactly $1$, so no point of the domain is left uncovered. A rectangle that has not been accepted at depth $22$ aborts the run.

The second program was written separately, does not import the first, and does not use the symmetry. It covers the full domain $[1/3,3]\times[0,1/2]$, of area $4/3$, by four-way subdivision, decides the sign of every completion in \eqref{eq:D1} by its own rationalisation of the numerator $2p_{1k}(rw_k\beta_{1,k}-2/(3w_k))$ rather than by a reused comparison, and carries a separate integrity calculation for the ordered tangent intersections. It terminates at depth $11$ with $466$ leaves whose areas sum to exactly $4/3$, with an abort at depth $12$.

\emph{Arithmetic.} Neither program uses floating point. All quantities are Python fractions or integers scaled by $Q=10^{12}$, and every inequality is evaluated with explicit and checked rounding directions: positive lower bounds are rounded downward, subtracted costs and denominator upper bounds upward. The tangent intercepts of \eqref{eq:tangent} are rounded downward to denominator $10^8$ and verified to be nonpositive and to satisfy $L_{\xi,m}+C_{\xi,m}\le1$, as an admissible lower bound must. The constant $\pi$ is enclosed by Machin's formula $16\arctan(1/5)-4\arctan(1/239)$ with alternating rational Leibniz remainders, the two programs using different series lengths $(14,6)$ and $(12,8)$, and the arctangent is enclosed by different recursively reduced series. The two decimal enclosures of Theorem~\ref{thm:main} are checked by comparing squared rational endpoints against the exact squared numerators,
\[
(2\pi_{\rm hi}c_{\rm lo})^2<\frac{341}{250}<(2\pi_{\rm lo}c_{\rm hi})^2,\qquad
(2\pi_{\rm hi}c_{\rm lo}')^2<\frac{6821}{5000}<(2\pi_{\rm lo}c_{\rm hi}')^2,
\]
so the printed digits follow from the rational certificate and not from a library constant.

\emph{What the certificate covers.} It covers the finite rational inequalities: the tail row costs \eqref{eq:qred} and the cap \eqref{eq:qcap}, the completion formulas \eqref{eq:D1} and \eqref{eq:D0} with their signs, the ordered breakpoints of the head envelope, the two covering trees with their areas and depth caps, and the three exterior comparisons including the tail sum below $7/75$. It also verifies by SHA-256 hash thirteen inherited proof and arithmetic files on which the argument rests, so that an accidental change to an inherited input cannot pass silently.

\emph{What the certificate does not cover.} It does not cover the Harnack inequalities \eqref{eq:harnack}, the width cone \eqref{eq:cone}, the Poisson envelope of Proposition~\ref{prop:poisson}, the positive-resolvent reduction of Section~\ref{sec:reduction}, or the derivation of the exterior formulas. Those are analytic statements, proved above in prose, and the programs only re-evaluate the rational comparisons that the prose leaves to arithmetic. The two programs also share the same decomposition parameters, $M=20$, the $7/10$ split, the eight tangent points and the $43$ head intervals, and they copy the same exterior block; this is an honest limit on their independence, and it means that the certificate tests the arithmetic of the covering rather than the choice of decomposition. That choice is not certified by computation; it is fixed by the estimates of Section~\ref{sec:proof}, each of which is an explicit inequality in the text.

\emph{Guarded failure.} Each program carries six deliberately corrupted runs, chosen so that each corrupts a different ingredient. Four corrupt the shared scaffold: a wrong pin hash, a wrong strict reserve ratio, a covering domain shrunk from $r\le3$ to $r\le2$, and a reserve raised past feasibility. The other two differ between the programs: the first alters the infinite-tail cost and replaces the centered tangent slope by one that is too steep to be an admissible lower bound, and the second corrupts the sign of a Stage-0 completion and gives a tangent intercept the wrong sign. Every one of the twelve must exit with status $1$ and a line beginning \texttt{REJECT}, rather than crash from a missing dependency, and the driver compares the observed status with the expected one and aborts otherwise. The purpose is narrow and worth stating: it shows that the acceptance predicates actually read the quantities they claim to read, and that a plausible arithmetic slip does not survive.

\section{Comparison, limitations and the open question}\label{sec:compare}

\begin{table}[t]
\centering
\small
\setlength{\tabcolsep}{4pt}
\begin{tabular}{@{}lllcc@{}}
\hline
constant & class and status & exact form & decimal & above $\Cinf$ \\
\hline
$\Cinf$ & scalar, exact, from~\cite{finite} & transcendental & $0.1827262478$ & \\
$\Cstar$ & centered, supremum, proved here & $\sqrt{341/250}/(2\pi)$ & $0.1858777115$ & $1.725\%$ \\
$\Crep$ & centered, strict, proved here & $\sqrt{6821/5000}/(2\pi)$ & $0.1858913384$ & $1.732\%$ \\
$\sqrt{683/500}/(2\pi)$ & centered, supremum, from~\cite{finiteenergy} & $\sqrt{683/500}/(2\pi)$ & $0.1860139355$ & $1.799\%$ \\
$\Ccent$ & centered, strict, from~\cite{finiteenergy} & $\sqrt{41/30}/(2\pi)$ & $0.1860593214$ & $1.824\%$ \\
$\Ccaus$ & causal, strict, from~\cite{finiteenergy} & $\sqrt{45/32}/(2\pi)$ & $0.1887345454$ & $3.288\%$ \\
$C_{21}$ & causal, single stage, from~\cite{optimum} & $\sqrt{377/250}/(2\pi)$ & $0.1954433070$ & $6.960\%$ \\
$C_1$ & causal, multipoint, from~\cite{optimum} & $\sqrt{413/250}/(2\pi)$ & $0.2045620904$ & $11.95\%$ \\
$U$ & centered, from~\cite{predecessor} & $\sqrt{6466219/3576050}/(2\pi)$ & $0.2140146706$ & $17.12\%$ \\
\hline
\end{tabular}
\caption{The exact scalar constant, the two constants proved here, and the successive coupled-channel ceilings. The decimals are rounded to ten places, and the last column is the relative excess over $\Cinf$. The causal constant is unchanged by this paper.}\label{tab:constants}
\end{table}

Table~\ref{tab:constants} lists the numbers. The ordering
\[
\Cinf<\Cstar<\Crep<\frac{\sqrt{683/500}}{2\pi}<\Ccent<\Ccaus<C_{21}<C_1<U
\]
is a chain of strict inequalities between explicit numbers, the middle three being the constants of~\cite{finiteenergy} and the last three those of~\cite{optimum,predecessor}. The new constant improves both constants of \cite{finiteenergy} for the centered class and touches nothing else. Measured as relative excess over $\Cinf$, the centered figure moves from $1.824\%$ to $1.732\%$ for the strict statement and from $1.799\%$ to $1.725\%$ for the supremum, so $4.1\%$ of the excess that was still open above $\Cinf$ has been closed for the supremum and $5.0\%$ for the strict statement. Measured as the fraction of the excess of $C_{21}$ that a constant removes, the figure moves from $73.79\%$ for the strict constant of~\cite{finiteenergy} to $75.11\%$ for $\Crep$ here.

\subsection{The floor of the method}

The edge-reserve argument has a limit that was proved in~\cite{finiteenergy}. Let $\mathcal W_{\rm cent}$ be the set of weight sequences of the centered class. An inequality for one edge of the form \eqref{eq:edge}, required to hold for every member of $\mathcal W_{\rm cent}$ at every $t>0$ and every $r>0$ with a reserve $d$ independent of $t$, $r$ and the weights, cannot have $d$ larger than $1-\Lambda/2$ where $\Lambda$ is the largest value of the ratio $-I_{xy}/B_r$ over that set, since \eqref{eq:edge} reads $-I_{xy}/B_r\le2(1-d)$. Explicit weight sequences and explicit trial vectors give $\Lambda>1.357186$, hence $d<0.321407$ and no such proof gives a constant below
\[
\frac{\sqrt{\Lambda}}{2\pi}\ >\ 0.185411 .
\]
This is a statement about the method, not a lower bound for $S_{\mathrm{cent}}$, and a proof that used the normal equation to tie the weights to the masses would not be subject to it. The reserve used here is $\dstar=0.318$, which is below the floor value $0.321407$ and therefore consistent with it, and the constant proved here exceeds the floor by $0.25\%$, so the floor lies at the fraction $0.852$ of $[\Cinf,\Cstar]$. The remaining $15\%$ of the interval is what a different proof, not a better implementation, would have to reach.

\subsection{What the result does and does not say}

The result is an upper bound on a linear onset temperature of a model class. Four limits should be kept in view.

First, the bound is linear. It is derived from the folded stability form of linearised Eliashberg theory, with the pairing kernel tested against a fixed positive normal branch. It makes no statement about the nonlinear gap equation at low temperature, and $\Tlin$ is not $\mathrm{T_c}$ for any exact Hamiltonian.

Second, the class is a subclass. The glue is retarded and the same nonnegative transfer enters the normal and pairing equations; vertex corrections, further anisotropy, magnetic and competing orders are outside, and $\Cinf$ is a property of this class, not a bound for any material.

Third, the widths used satisfy the cone \eqref{eq:cone}. A branch whose widths leave the cone is not covered, and the causal class, where only \eqref{eq:harnack} is retained, keeps the larger constant $\Ccaus$ of~\cite{finiteenergy}.

Fourth, the two bounds are different statements and both are needed. The supremum bound $S_{\mathrm{cent}}\le\Cstar$ gives no strict inequality at a particular model, since endpoint stability without a neighbourhood theorem would not justify it; the strict statement follows from it through $\Cstar<\Crep$, the bare gap $1/6821$ recording the slack at the reported temperature. Neither requires monotonicity of the onset in temperature.

\subsection{The open question}

Let $S_{\mathrm{cent}}$ be the supremum defined in Theorem~\ref{thm:main}. From \cite{finite,predecessor} one has $\Cinf\le S_{\mathrm{cent}}$, because the scalar case belongs to the class and its onset approaches $\Cinf\sqrt\rho$ along soft sequences, and from Theorem~\ref{thm:main} one has $S_{\mathrm{cent}}\le\Cstar$. Thus
\[
\Cinf\ \le\ S_{\mathrm{cent}}\ \le\ \Cstar,\qquad \frac{\Cstar}{\Cinf}-1=1.725\% .
\]
Whether $S_{\mathrm{cent}}=\Cinf$ or $S_{\mathrm{cent}}>\Cinf$ is open. Nothing in this paper or in \cite{predecessor,optimum,finiteenergy} decides it. For a single channel the supremum is $\Cinf$, and for coupled channels with energy-independent density of states the bound $\Cinf$ holds as well, so the question concerns only several channels with different energy-dependent profiles. The exploratory float searches recorded in~\cite{finiteenergy} found no witness above $\Cinf$, but they enclose neither the finite pairing blocks nor the truncated normal tails, so they prove nothing either way. The two natural routes to a smaller constant are a comparison that retains the coupled normal masses across channels and frequencies at once, and an arrangement of the operator that leaves the edge-by-edge form; the floor excludes neither. A certified excess witness, if one exists, would have to be an actual branch rather than a frozen weight sequence.

\section{Reproducibility and data availability}\label{sec:repro}

A replay supplement accompanies this preprint. It contains the written proof from which this paper is drawn, the two exact covering programs, the mutation driver, the closure pin file with thirteen SHA-256 hashes of inherited proofs and arithmetic, and the execution transcript. The programs use only the Python standard library; no interval package, no compiled extension and no network access are involved.

From the directory holding the supplement one runs
\begin{center}\texttt{python -B verify.py}\end{center}
and obtains exit status $0$ when everything passes. The driver runs the primary cover, which must exit $0$ and print the certified covering with $159$ leaves, maximum depth $20$ and area $1$; the independent cover, which must exit $0$ and print $466$ leaves, maximum depth $11$ and area $4/3$; and twelve deliberately corrupted runs, six of each program, each of which must exit $1$ and print a line beginning \texttt{REJECT}. Both base runs also print the exact decimal enclosures of the two constants. The whole run takes about two and a half minutes. Everything analytic, namely Proposition~\ref{prop:poisson}, the Harnack inequalities, the width cone, the reduction of Section~\ref{sec:reduction} and the derivation of the exterior estimates, is in the proof above and not in the replay.

\emph{Data availability.} This work uses no experimental data and no material samples. All computational evidence is the replay supplement, released as DOI \DOI{\SUPPLEMENTDOI}, which holds the two checkers, the mutation driver, the pin file and the transcript; the paper record is DOI \DOI{\PAPERDOI}. The earlier results on which the argument rests are cited by DOI in the references.

\section*{Supporting material and disclosure}

This preprint has not been peer reviewed. The result is a theorem about linearised Eliashberg theory. It is an upper bound on a linear model onset temperature, not an exact transition temperature and not a measurement, and no material, synthesis, prediction of a compound or claim that any substance is superconducting is made anywhere in it. AI agents generated the proofs, code, and manuscript drafts under D.~Rens' direction. Separate AI agents performed the independent reviews. One review covered the written proof and its reduction, and a second covered the certificates and their execution record. The replay supplement accompanies the preprint.

\begingroup
\raggedright
\bibliographystyle{plain}
\bibliography{refs}
\endgroup

\end{document}
